\makeabstract
    {
    Investigating drying behavior in colloid{\textendash}polymer mixtures using multiparticle collision dynamics 
    }
    {
    Samaira Abdul-Waalee\textsuperscript{1}, Jinny Cha\textsuperscript{2}, Aanuoluwa Amusan\textsuperscript{2}, Michael P. Howard\textsuperscript{2}
    }
    {
    \textsuperscript{1} Amherst College, Amherst, MA \\
\textsuperscript{2} Department of Chemical Engineering, Auburn University, Auburn, AL
    }
    {
    Colloid{\textendash}polymer mixtures are used in coatings, printed electronics, and battery electrodes, where drying affects material structure and performance. This project investigated how polymer chain length and particle interactions influence drying behavior using Multiparticle Collision Dynamics (MPCD) simulations.
    }
    {
    Multiparticle Collision Dynamics (MPCD) simulations were performed using Python on the Delta computing system to model a drying film containing colloids, polymers, and solvent. Unlike Brownian Dynamics, MPCD captures solvent flow and hydrodynamic interactions. Simulations were performed at constant polymer (0.005) and colloid (0.1) concentrations with polymer chain lengths of 10, 20, and 40. Attractive and non-attractive interactions were compared to examine their effects on drying behavior.
    }
    {
    
As drying progressed, colloids moved toward the drying surface and became more concentrated at the interface. Polymers also moved during drying, becoming more concentrated in the film's interior rather than at the surface. Polymer chain lengths of 10, 20, and 40 showed similar colloid{\textendash}polymer distributions throughout drying, indicating that chain length had only a minor effect on the final film structure. The interaction type had a stronger effect than the chain length. Non-attractive interactions maintained a more uniform particle distribution, whereas attractive interactions led to particle clustering. These findings show that particle interactions strongly influence how colloids and polymers organize during drying and affect the final structure of the dried film.
    }
    {
    Similar to previous Brownian Dynamics (BD) studies, colloids accumulated near the drying surface during evaporation. However, our Multiparticle Collision Dynamics (MPCD) simulations showed a stronger accumulation of colloids at the interface. In contrast to previous BD studies that reported a sharp polymer peak near the drying surface, our simulations showed a broader and more uniform polymer distribution. Polymer chain length had only a small effect on particle distribution. Attractive interactions caused particles to cluster, while non-attractive interactions maintained a more even particle distribution. These results show that including solvent flow and particle interactions changes how colloids and polymers redistribute during drying and influences the final structure of the dried film.
    }

    \newpage
    